% Generated draft: requires mathematical and source review.
\documentclass[11pt]{article}
\usepackage[margin=25mm]{geometry}
\usepackage{fontspec}
\setmainfont{Latin Modern Roman}
\usepackage{amsmath,amssymb,mathtools,unicode-math}
\setmathfont{Latin Modern Math}
\usepackage{enumitem,needspace,xurl,hyperref}
\setlength{\parindent}{0pt}
\setlength{\parskip}{5pt}
\setlength{\emergencystretch}{3em}
\newcommand{\N}{\mathbb{N}}
\newcommand{\Z}{\mathbb{Z}}
\newcommand{\Q}{\mathbb{Q}}
\newcommand{\R}{\mathbb{R}}
\newcommand{\C}{\mathbb{C}}
\newcommand{\F}{\mathbb{F}}
\newcommand{\A}{\mathbb{A}}
\newcommand{\PP}{\mathbb{P}}
\newcommand{\E}{\mathbb{E}}
\newcommand{\eps}{\varepsilon}
\newcommand{\dd}{\,\mathrm{d}}
\begin{document}
\section*{UnsolvedMath ID: 20000344}
\textbf{Source problem code:} catalogue entry
\subsection{Definitions and mathematical statement}
Let $\overline M_g$ be the projective coarse moduli space of stable genus-$g$ curves, interpreted over the intended characteristic-zero field, usually $\C$. Its Kodaira dimension means the birational Kodaira dimension of a smooth projective model, not the logarithmic Kodaira dimension of the open smooth-curve locus. The source writes $\overline{\mathcal M}_g$; coarse-space versus stack conventions should be checked.

\textbf{Original Problem 5.2.} Determine
\[\kappa(\overline M_g)\qquad\text{for }g=15\text{ or }17\le g\le22.\]
The original historical background says the spaces are of general type for $g>22$ and that the other then-known cases have negative Kodaira dimension. The supplied verification restores the bars lost by extraction and removes the following subsection heading ``5.2. Other questions'' from the problem.

\textbf{Attached unreviewed status account.} The aid reports $\kappa(\overline M_{15})=-\infty$, unknown Kodaira dimension for $17\le g\le21$, and general type for $g=22$ (hence Kodaira dimension $3g-3$ there). These are the supplied record's genus-by-genus literature claims, not a newly verified current-status assessment; no supporting publications are identified.

The aid also reports a surface-pencil criterion. Let $S$ be a smooth projective surface, $L$ a line bundle, and take an irreducible-nodal Lefschetz pencil in $|L|$ whose smooth fibers have genus $g\ge2$ and whose only stable degenerations are irreducible one-node curves. Blow up its base points to obtain a stable family and a moduli test curve $B$ with its family multiplicity convention understood. The asserted coarse-canonical intersection formula is
\[2K_{\overline M_g}\cdot B=4K_S^2-22\chi(\mathcal O_S)+L^2+5K_S\cdot L.\]
The qualifier ``meeting only $\delta_0$'' excludes separating-node boundary contributions. Exact assumptions on the pencil, automorphisms, stable reduction, coarse intersection and degree of its classifying map are not supplied and must be reviewed.

For a K3 pencil, the reported consequence is
\[K_{\overline M_g}\cdot B=g-23.\]
In the open genus range $17\le g\le21$ this is negative, but the aid says the K3-section construction has parameter dimension at most $19+g$, and therefore codimension at least $2g-22$ in dimension $3g-3$. It is not covering. The BDPP implication invoked by the aid requires a genuinely movable/covering curve realization on the appropriate projective birational model; negative intersection with these confined K3 pencils alone does not establish non-pseudoeffectivity of the canonical class or negative Kodaira dimension of the whole moduli space.

The aid additionally mentions an exact standard-divisor wall ledger, a positive-cone obstruction, and an explicit slope inequality for effective divisors avoiding a proposed surface family. Their divisor coefficients, inequalities and hypotheses are absent from the supplied summary, so they are recorded as unsupplied parts of the research attempt rather than reconstructed or asserted here.
\subsection{Short English statement}
The problem concerns the compactified stable-curve moduli spaces in specified genera. The aid reports the genus-by-genus status and a canonical-intersection formula for surface pencils. Its negative K3-pencil intersections occur in a locus of positive codimension, so the reported criterion requires an additional covering-family condition.
\subsection{Sources}
\begin{itemize}
\item Local UnsolvedMath registry, permanent ID 20000344, source code catalogue entry; \texttt{aim-algebraic-geometry-notes.json}, notes[343], tag \texttt{problem}.
\item AIM workshop \emph{Compact moduli spaces and birational geometry}, Problem 5.2: \url{https://aimath.org/WWN/birational/birational.pdf}. The supplied \texttt{corrected\_verified} field explains the restored compactification bars and excluded following heading. No PDF inspection is claimed here.
\item Attached research attempt 1 gives the modern-status account and mentions BDPP without bibliographic locators; original literature field empty. Result \texttt{partial\_result}, status \texttt{partially\_solved}, rigor label \texttt{proved}, novelty confidence low. It is machine generated, not peer reviewed, and requires independent verification before citation.
\end{itemize}
\end{document}
